% source: manti.01.ma, 1ère SM, devoir (Drive 0B6pfSJcEVng0d1ZURHVDZllJa00). Ex.2 logique (hors thème)
\ex{Fonctions}{d1ZU}{Ex. 1}{7}
On considère la fonction $f$ définie par : $f(x)=\dfrac{2x}{x+1}$ si $x\leq1$ \ et \ $f(x)=\dfrac{2x+1}{x+2}$ si $x>1$.
\begin{enumerate}
\item Déterminer l'ensemble de définition de $f$.
\item Soit $h$ la fonction définie par $h(x)=\dfrac2x$.
\begin{enumerate}
\item Résoudre l'inéquation $h(x)>1$.
\item Montrer que l'ensemble de définition de $f\circ h$ est $D=\mathbb{R}\setminus\{-2,0\}$.
\item Calculer $(f\circ h)(x)$ pour tout $x$ de $D$.
\end{enumerate}
\end{enumerate}
\fin
\ex{Fonctions}{d1ZU}{Ex. 3}{9}
On considère la fonction $h$ définie sur $[0,+\infty[$ par $h(x)=x^3-3x+2$.
\begin{enumerate}
\item \begin{enumerate}
\item Montrer que $T_h(x,y)=x^2+y^2+xy-3$.
\item Étudier la monotonie de $h$ sur chacun des intervalles $[0,1]$ et $[1,+\infty[$.
\end{enumerate}
\item Montrer que si $a_1,a_2,\dots,a_n$ sont des réels strictement positifs tels que $a_1\times a_2\times\dots\times a_n=1$, alors $\left(2+a_1^3\right)\left(2+a_2^3\right)\cdots\left(2+a_n^3\right)\geq3^n$.
\end{enumerate}
\fin
